\section{Related Work}
\label{sec:related}

\subsection{Data Association in Multi-Object Tracking}

Data association in \MOT is a classical linear assignment problem. 
Traditional methods include:
\begin{itemize}[leftmargin=*]
    \item \textbf{Hungarian algorithm}~\cite{kuhn1955hungarian}: 
          an $O(n^3)$ exact polynomial algorithm.
    \item \textbf{MIP methods}~\cite{jiang2015multiple}: integer 
          programming, exact at small scales but slow at large scales.
    \item \textbf{Greedy matching}: $O(n \log n)$ but non-optimal.
    \item \textbf{Deep learning approaches}~\cite{zhang2022bytetrack,
          cao2023observation}: learn similarity functions, but 
          association itself still uses Hungarian.
\end{itemize}

A key observation is that the \emph{combinatorial} structure of \MOT 
data association is \emph{linear}. Its dual form admits exact polynomial 
solution. This is the theoretical reason for Hungarian's dominance.

\subsection{Cooperative Perception}

Cooperative perception~\cite{chen2019cooper} enables multiple vehicles 
or edge devices to share sensor data and overcome single-vehicle field-
of-view limitations. Representative methods include:
\begin{itemize}[leftmargin=*]
    \item \textbf{V2X-ViT}~\cite{xu2022v2x}: Transformer-based 
          cooperative perception.
    \item \textbf{CoBEVT}~\cite{xu2022cobevt}: bird's-eye view 
          cooperative perception.
    \item \textbf{AttFuse}~\cite{xu2022opv2v}: attention-based 
          intermediate fusion.
\end{itemize}
Again, the association step in these pipelines remains Hungarian.

\subsection{Physics-Native Computing and \abnqss}

\abnqss v3.0~\cite{chang2026abnqss} is a representative 
physics-native computing architecture with:
\begin{itemize}[leftmargin=*]
    \item \textbf{L0 photonic layer}: silicon-photonics MZI 
          interferometer array for matrix multiplication and 
          null-space projection.
    \item \textbf{L1 analog layer}: OTA-C ``Shenqu engine'' for 
          nonlinear feedback control.
    \item \textbf{L2 digital layer}: FPGA/CMOS for timing and 
          buffering.
    \item \textbf{L5 cognitive layer}: external host/SoC for 
          macro task scheduling.
\end{itemize}
The digital twin of \abnqss v3.0 has been validated for $N=3$ to 
$N=6$ via Python simulation~\cite{chang2026n4}.

\subsection{Hungarian Algorithm and the Dual Structure}

Linear assignment is:
\begin{equation}
    \min_{x \in \{0,1\}^{T \times D}} \sum_{i,j} c_{ij} x_{ij}
    \quad \text{s.t.} \quad 
    \sum_i x_{ij} = 1,\; \sum_j x_{ij} = 1.
\end{equation}
Its Lagrangian dual is:
\begin{equation}
    \max_{u, v} \sum_i u_i + \sum_j v_j
    \quad \text{s.t.} \quad u_i + v_j \le c_{ij}.
\end{equation}
The Hungarian algorithm exploits this dual structure and terminates 
in $O(n^3)$ time with an exact certificate of optimality.

\textbf{Theoretical consequence}: any iterative approximation 
algorithm (including \ZSP) cannot improve on Hungarian for this 
problem class, because the dual certificate is already exact.